% GENERATED FILE. Edit canonical lessons, then run npm run generate:books.
% canonical-source-hash: fnv1a64:cd21ba18c4713a1f
% canonical-lessons: TA-C421-valikatti, TA-C421-kaiyetu, TA-C421-currikkattu, TA-C421-parvaiyitu, TA-C421-valiyanuppu

\chapter{Guides and Guidebooks}
\label{ch:ta-guides-and-guidebooks}
\hlchaptermodality{421}

\begin{chapteropening}
\textbf{By the end of this chapter:} \emph{I can talk about guides and guidebooks, showing someone around, visiting and seeing people off.}

Complete the last lesson of chapter 421: I can talk about guides and guidebooks, showing someone around, visiting and seeing people off.
\end{chapteropening}

\section[\texorpdfstring{\ta{வழிகாட்டி} (vaḻikāṭṭi)}{வழிகாட்டி (vaḻikāṭṭi)}]{\ta{வழிகாட்டி} (vaḻikāṭṭi) — a guide}

\label{lesson:TA-C421-valikatti}



\begin{quote}
Before the new one: say the Tamil for a wage, then the Tamil for cleaning.
\end{quote}

\subsection*{You'll want to know: \ta{வழிகாட்டி}}
\textbf{\ta{வழிகாட்டி}} — \emph{vaḻikāṭṭi} — \textquotedblleft{}a guide\textquotedblright{}.

\begin{quote}
\ta{வழிகாட்டி} \ta{அரண்மனையின்} \ta{கதையைச்} \ta{சொன்னார்}.
\end{quote}

\begin{quote}
\emph{vaḻikāṭṭi araṇmaṉaiyiṉ kataiyaic coṉṉār.}
\end{quote}

\begin{quote}
The guide told the story of the palace.
\end{quote}

\ta{வழி}, the way, + \ta{காட்டி}, one who shows.

\begin{cousinweb}
\textbf{\ta{வழிகாட்டி}} (\emph{vaḻikāṭṭi}) is built from \ta{வழி}, \textquotedblleft{}a way\textquotedblright{} (DEDR 5297), and \ta{காட்டி}, \textquotedblleft{}one who shows\textquotedblright{}, from \ta{காட்டு} (DEDR 1443).
\end{cousinweb}

\begin{grammarlens}[title={What you've built}]
One more word for getting around.
\end{grammarlens}

\subsection*{Your turn}
\begin{itemize}
\raggedright
  \item \emph{Say it:} \emph{vaḻikāṭṭi}
  \item \emph{Say it:} \emph{vaḻikāṭṭi}, once more
  \item \emph{Say it:} \emph{vaḻikāṭṭi araṇmaṉaiyiṉ kataiyaic coṉṉār.} — The guide told the story of the palace.
  \item \emph{Recall:} say the Tamil for a wage, then the Tamil for cleaning, then say \emph{vaḻikāṭṭi} again
\end{itemize}

\begin{tcolorbox}[breakable,colback=teal!4,colframe=teal!35!black,title={Before you move on}]
What does \ta{வழிகாட்டி} mean? (\textquotedblleft{}A guide\textquotedblright{}.) Say it once more.
\end{tcolorbox}
\section[\texorpdfstring{\ta{கையேடு} (kaiyēṭu)}{கையேடு (kaiyēṭu)}]{\ta{கையேடு} (kaiyēṭu) — a handbook; a guidebook}

\label{lesson:TA-C421-kaiyetu}



\begin{quote}
Before the new one: say the Tamil for cleaning, then the Tamil for a guide.
\end{quote}

\subsection*{You'll want to know: \ta{கையேடு}}
\textbf{\ta{கையேடு}} — \emph{kaiyēṭu} — \textquotedblleft{}a handbook; a guidebook\textquotedblright{}.

\begin{quote}
\ta{அந்தக்} \ta{கையேட்டில்} \ta{நல்ல} \ta{உணவகங்களின்} \ta{பட்டியல்} \ta{இருக்கிறது}.
\end{quote}

\begin{quote}
\emph{antak kaiyēṭṭil nalla uṇavakaṅkaḷiṉ paṭṭiyal irukkiṟatu.}
\end{quote}

\begin{quote}
That guidebook has a list of good restaurants.
\end{quote}

A -\ta{டு} noun doubles before a case: \ta{கையேட்டில்}, in the handbook.

\begin{cousinweb}
\textbf{\ta{கையேடு}} (\emph{kaiyēṭu}) is built from \ta{கை}, \textquotedblleft{}hand\textquotedblright{} (DEDR 2023), and \ta{ஏடு}, \textquotedblleft{}a palm leaf for writing, a book\textquotedblright{} (DEDR 884): a book for the hand.
\end{cousinweb}

\begin{grammarlens}[title={What you've built}]
One more word for getting around.
\end{grammarlens}

\subsection*{Your turn}
\begin{itemize}
\raggedright
  \item \emph{Say it:} \emph{kaiyēṭu}
  \item \emph{Say it:} \emph{kaiyēṭu}, once more
  \item \emph{Say it:} \emph{antak kaiyēṭṭil nalla uṇavakaṅkaḷiṉ paṭṭiyal irukkiṟatu.} — That guidebook has a list of good restaurants.
  \item \emph{Recall:} say the Tamil for cleaning, then the Tamil for a guide, then say \emph{kaiyēṭu} again
\end{itemize}

\begin{tcolorbox}[breakable,colback=teal!4,colframe=teal!35!black,title={Before you move on}]
What does \ta{கையேடு} mean? (\textquotedblleft{}A handbook; a guidebook\textquotedblright{}.) Say it once more.
\end{tcolorbox}
\section[\texorpdfstring{\ta{சுற்றிக்காட்டு} (cuṟṟikkāṭṭu)}{சுற்றிக்காட்டு (cuṟṟikkāṭṭu)}]{\ta{சுற்றிக்காட்டு} (cuṟṟikkāṭṭu) — to show around}

\label{lesson:TA-C421-currikkattu}



\begin{quote}
Before the new one: say the Tamil for a guide, then the Tamil for a handbook.
\end{quote}

\subsection*{You'll want to know: \ta{சுற்றிக்காட்டு}}
\textbf{\ta{சுற்றிக்காட்டு}} — \emph{cuṟṟikkāṭṭu} — \textquotedblleft{}to show around\textquotedblright{}.

\begin{quote}
\ta{என்} \ta{நண்பன்} \ta{எனக்கு} \ta{ஊரைச்} \ta{சுற்றிக்காட்டினான்}.
\end{quote}

\begin{quote}
\emph{eṉ naṇpaṉ eṉakku ūraic cuṟṟikkāṭṭiṉāṉ.}
\end{quote}

\begin{quote}
My friend showed me around the town.
\end{quote}

Past (I …): \ta{சுற்றிக்காட்டினேன்}, \emph{cuṟṟikkāṭṭiṉēṉ}.

The person shown takes -\ta{க்கு}, the place -\ta{ஐ}: \ta{எனக்கு} \ta{ஊரைச்} \ta{சுற்றிக்காட்டு}.

\begin{cousinweb}
\textbf{\ta{சுற்றிக்காட்டு}} (\emph{cuṟṟikkāṭṭu}) is built from \ta{சுற்றி}, \textquotedblleft{}going around\textquotedblright{}, from \ta{சுற்று} (DEDR 2715), and \ta{காட்டு}, \textquotedblleft{}to show\textquotedblright{} (DEDR 1443).
\end{cousinweb}

\begin{grammarlens}[title={What you've built}]
One more word for getting around.
\end{grammarlens}

\subsection*{Your turn}
\begin{itemize}
\raggedright
  \item \emph{Say it:} \emph{cuṟṟikkāṭṭu}
  \item \emph{Say it:} \emph{cuṟṟikkāṭṭu}, once more
  \item \emph{Say it:} \emph{eṉ naṇpaṉ eṉakku ūraic cuṟṟikkāṭṭiṉāṉ.} — My friend showed me around the town.
  \item \emph{Recall:} say the Tamil for a guide, then the Tamil for a handbook, then say \emph{cuṟṟikkāṭṭu} again
\end{itemize}

\begin{tcolorbox}[breakable,colback=teal!4,colframe=teal!35!black,title={Before you move on}]
What does \ta{சுற்றிக்காட்டு} mean? (\textquotedblleft{}To show around\textquotedblright{}.) Say it once more.
\end{tcolorbox}
\section[\texorpdfstring{\ta{பார்வையிடு} (pārvaiyiṭu)}{பார்வையிடு (pārvaiyiṭu)}]{\ta{பார்வையிடு} (pārvaiyiṭu) — to visit; to inspect}

\label{lesson:TA-C421-parvaiyitu}



\begin{quote}
Before the new one: say the Tamil for a handbook, then the Tamil for to show around.
\end{quote}

\subsection*{You'll want to know: \ta{பார்வையிடு}}
\textbf{\ta{பார்வையிடு}} — \emph{pārvaiyiṭu} — \textquotedblleft{}to visit; to inspect\textquotedblright{}.

\begin{quote}
\ta{சுற்றுலாப்} \ta{பயணிகள்} \ta{கோட்டையைப்} \ta{பார்வையிட்டார்கள்}.
\end{quote}

\begin{quote}
\emph{cuṟṟulāp payaṇikaḷ kōṭṭaiyaip pārvaiyiṭṭārkaḷ.}
\end{quote}

\begin{quote}
The tourists visited the fort.
\end{quote}

Past (I …): \ta{பார்வையிட்டேன்}, \emph{pārvaiyiṭṭēṉ}.

Formal: of officials and visitors; in speech, \ta{போய்ப்} \ta{பார்}.

\begin{cousinweb}
\textbf{\ta{பார்வையிடு}} (\emph{pārvaiyiṭu}) is built from \ta{பார்வை}, \textquotedblleft{}a look, sight\textquotedblright{} (DEDR 4091), and \ta{இடு}, \textquotedblleft{}to put, to give\textquotedblright{} (DEDR 442).
\end{cousinweb}

\begin{grammarlens}[title={What you've built}]
One more word for getting around.
\end{grammarlens}

\subsection*{Your turn}
\begin{itemize}
\raggedright
  \item \emph{Say it:} \emph{pārvaiyiṭu}
  \item \emph{Say it:} \emph{pārvaiyiṭu}, once more
  \item \emph{Say it:} \emph{cuṟṟulāp payaṇikaḷ kōṭṭaiyaip pārvaiyiṭṭārkaḷ.} — The tourists visited the fort.
  \item \emph{Recall:} say the Tamil for a handbook, then the Tamil for to show around, then say \emph{pārvaiyiṭu} again
\end{itemize}

\begin{tcolorbox}[breakable,colback=teal!4,colframe=teal!35!black,title={Before you move on}]
What does \ta{பார்வையிடு} mean? (\textquotedblleft{}To visit; to inspect\textquotedblright{}.) Say it once more.
\end{tcolorbox}
\section[\texorpdfstring{\ta{வழியனுப்பு} (vaḻiyaṉuppu)}{வழியனுப்பு (vaḻiyaṉuppu)}]{\ta{வழியனுப்பு} (vaḻiyaṉuppu) — to see off}

\label{lesson:TA-C421-valiyanuppu}



\begin{quote}
Before the new one: say the Tamil for to show around, then the Tamil for to visit.
\end{quote}

\subsection*{You'll want to know: \ta{வழியனுப்பு}}
\textbf{\ta{வழியனுப்பு}} — \emph{vaḻiyaṉuppu} — \textquotedblleft{}to see off\textquotedblright{}.

\begin{quote}
\ta{நாங்கள்} \ta{தாத்தாவை} \ta{ரயில்} \ta{நிலையத்தில்} \ta{வழியனுப்பினோம்}.
\end{quote}

\begin{quote}
\emph{nāṅkaḷ tāttāvai rayil nilaiyattil vaḻiyaṉuppiṉōm.}
\end{quote}

\begin{quote}
We saw grandfather off at the railway station.
\end{quote}

Past (I …): \ta{வழியனுப்பினேன்}, \emph{vaḻiyaṉuppiṉēṉ}.

The person seen off takes -\ta{ஐ}: \ta{தாத்தாவை} \ta{வழியனுப்பு}.

\begin{cousinweb}
\textbf{\ta{வழியனுப்பு}} (\emph{vaḻiyaṉuppu}) is built from \ta{வழி}, \textquotedblleft{}a way\textquotedblright{} (DEDR 5297), and \ta{அனுப்பு}, \textquotedblleft{}to send; to go a little way with someone\textquotedblright{} (DEDR 329).
\end{cousinweb}

\begin{grammarlens}[title={What you've built}]
One more word for getting around.
\end{grammarlens}

\subsection*{Your turn}
\begin{itemize}
\raggedright
  \item \emph{Say it:} \emph{vaḻiyaṉuppu}
  \item \emph{Say it:} \emph{vaḻiyaṉuppu}, once more
  \item \emph{Say it:} \emph{nāṅkaḷ tāttāvai rayil nilaiyattil vaḻiyaṉuppiṉōm.} — We saw grandfather off at the railway station.
  \item \emph{Recall:} say the Tamil for to show around, then the Tamil for to visit, then say \emph{vaḻiyaṉuppu} again
\end{itemize}

\begin{tcolorbox}[breakable,colback=teal!4,colframe=teal!35!black,title={Before you move on}]
What does \ta{வழியனுப்பு} mean? (\textquotedblleft{}To see off\textquotedblright{}.) Say it once more.
\end{tcolorbox}
